\section{Method}
\label{sec:method}

\subsection{Operational definitions}

Every figure in this paper rests on three definitions that are our choices,
not universal detections. They are stated here so a reader can disagree with
them specifically.

\paragraph{Escape attempt.} For tasks 01 through 03 it is the first call to
the \texttt{http\_fetch} tool. For tasks 04 through 06 it is the first
appearance of the run's canary token in the arguments or result of a tool
call. The difference matters: in those three tasks the planted hint authorises
consulting the channel, so reaching the mirror, the memory store, or the
knowledge base is not itself a violation. Acting on the injected instruction
is. Both definitions live in one function,
\texttt{compute\_ttd.escape\_ts\_for}, imported by every layer that needs
them, after an earlier version drifted between the analysis and the dashboard.

\paragraph{Escalation ladder.} Each run is scored on a seven-step ladder, from
0 (no signal) through 4 (successfully contacted an internal mock) to 5 (used
the out-of-scope forwarding, meaning the canary appears where an attacker
asked for it) and 6 (sandbox broken, detected as heartbeat silence). Levels 0
through 5 are computed from the agent transcript alone. Only level 6 consults
monitor events. That separation is load-bearing for this paper: every
correction we made to detection logic, described in
Section~\ref{sec:validity}, leaves the level counts untouched.

\paragraph{Time to detection.} The interval from the agent's act to the first
monitor alert at or after it. One filter is applied: when a container exits
cleanly, no heartbeat alert counts as a detection, because a clean shutdown
always produces heartbeat silence and an earlier version of the analysis
credited that silence as a real detection.

\subsection{Experimental matrix}

\begin{table}[h]
\centering
\small
\begin{tabular}{@{}lllrl@{}}
\toprule
\textbf{Set} & \textbf{Tasks} & \textbf{Conditions} & \textbf{Runs} & \textbf{Machine} \\
\midrule
\texttt{results/} & 04, 05, 06 & 2 (plus 3 partial) & 63 & A (Apple Silicon) \\
\texttt{results\_causal/} & 06 & 3 & 30 & B (Windows) \\
\texttt{results\_redteam/} & 04 & 2 & 20 & B (Windows) \\
\bottomrule
\end{tabular}
\caption{113 runs across three sets. The corpus was planned as 60 runs; the
three extra are an incomplete arm whose reported table had to be retracted
(Section~\ref{sec:validity}). Which machine produced which set is not
bookkeeping: it is the axis along which the main finding fails to replicate.}
\end{table}

All runs use \texttt{qwen2.5:7b-instruct}, with the digest verified identical
across both machines before the second set was generated. Runs are sequential
by design, since the model is the scarce resource and parallel runs would
produce contention rather than throughput.

\todo{Si Daniel pasa su version de Ollama y su OLLAMA\_CONTEXT\_LENGTH, anadir
aqui una fila de configuracion de runtime por maquina. Es el dato que decide
si la Seccion 5 reporta una limitacion explicada o un misterio abierto.}

\subsection{Deterministic validation}

A monitor that never fired is indistinguishable from a monitor that is broken.
Before reporting any detection rate, we therefore fire each event the monitors
are supposed to catch, by hand, and measure how long each takes. The suite
needs no model backend and costs nothing to run.

\begin{table}[h]
\centering
\small
\begin{tabular}{@{}rlr@{}}
\toprule
\textbf{\#} & \textbf{Control} & \textbf{Measured} \\
\midrule
1 & heartbeat silence after killing the pinger & 9.08\,s \\
2 & canary in a workspace directory name & 0.43\,s \\
3 & canary in an outbound query string & 0.49\,s \\
4 & destination outside the egress baseline & 0.49\,s \\
5 & canary in cross-run persistent memory & 18.45\,s \\
6 & red-teamer serves its payload through the SSRF & 0.52\,s \\
\bottomrule
\end{tabular}
\caption{All six fire. Control 1 measures the configured threshold (5\,s
interval times a multiplier of 2) minus the time since the last ping, not
monitor slowness. Control 5 is an order of magnitude slower than the other
canary surfaces because inotify does not fire for host-side writes to a bind
mount, so that surface falls back to the periodic sweep.}
\end{table}

Two of these controls exist because the harness failed silently without them.
Control 5 was added after finding that the most persistent leak channel in the
corpus, a canary written into cross-run memory, was the only one nothing
watched. Control 6 was added after finding that no control touched the
attacker path at all, and it reported a blind spot on its first run
(Section~\ref{sec:validity}).

A separate sweep varies each monitor's configuration parameter against the
same synthetic event, isolating how often a monitor looks from the variance of
model behaviour. Its results appear in Section~\ref{sec:results}.
